% Incorporar desde el documento grupal con % Incorporar desde el documento grupal con % Incorporar desde el documento grupal con % Incorporar desde el documento grupal con \input{notas/carlos.tex}
% si el archivo principal está en la raíz. Las citas usan referencias.bib.

\subsection{Carlos Manotoa · MQTT aplicado a CLAB}

\paragraph{Arquitectura y comunicación.}
CLAB gestiona los laboratorios informáticos de la UTEQ. Cuando el
administrador marca una falla como resuelta, Angular envía una petición REST a
Spring Boot y espera su respuesta. El backend actualiza el reporte en
PostgreSQL, guarda el seguimiento y avisa al docente que reportó la falla.
Para el administrador, todo empieza con un botón. La aplicación se encarga del
acceso a la base de datos \cite[sección~3.2 y caso CU-04]{clab2026}.

\paragraph{Despliegue y CAP/PACELC.}
El diagrama muestra un servidor PostgreSQL. Si Spring Boot pierde la conexión
con ese servidor, no hay un mecanismo de conmutación automática documentado.
Para aplicar CAP y PACELC haría falta conocer la replicación, el comportamiento
ante particiones y cómo se resuelve el compromiso entre latencia y consistencia
en operación normal. El informe permite identificar la dependencia de la base
de datos, pero deja esos puntos abiertos \cite[figura~4]{clab2026}.

\paragraph{Fallo observado durante las pruebas.}
En una prueba con Postman, un bloqueo por evento institucional dejó sin aviso
a los docentes con reservas afectadas. El equipo corrigió una condición en
\texttt{BloqueoService} durante el mismo sprint. Para verificar el flujo
completo, comprobaría que se registre el bloqueo y se genere el aviso \cite[sección~6.5.2]{clab2026}.

\paragraph{Propuesta y demostración.}
Para la resolución de fallas, propongo guardar el evento pendiente en la misma
transacción que actualiza el reporte. Así, si Mosquitto está caído, un proceso
puede reintentar la publicación y el aviso se puede demorar sin deshacer la
resolución. En MQTTX publico \texttt{ReporteFallaResuelto} en
\texttt{clab/fallas/resuelta} y muestro el mismo mensaje en Notificaciones y
Auditoría. La demo comprueba la entrega a dos suscriptores conectados con QoS~1.
El broker distribuye por tópico. Con QoS~1, la confirmación es por tramo y
pueden llegar duplicados. Cada consumidor podría reconocerlos mediante
\texttt{eventoId} \cite{mqtt2019}.

% si el archivo principal está en la raíz. Las citas usan referencias.bib.

\subsection{Carlos Manotoa · MQTT aplicado a CLAB}

\paragraph{Arquitectura y comunicación.}
CLAB gestiona los laboratorios informáticos de la UTEQ. Cuando el
administrador marca una falla como resuelta, Angular envía una petición REST a
Spring Boot y espera su respuesta. El backend actualiza el reporte en
PostgreSQL, guarda el seguimiento y avisa al docente que reportó la falla.
Para el administrador, todo empieza con un botón. La aplicación se encarga del
acceso a la base de datos \cite[sección~3.2 y caso CU-04]{clab2026}.

\paragraph{Despliegue y CAP/PACELC.}
El diagrama muestra un servidor PostgreSQL. Si Spring Boot pierde la conexión
con ese servidor, no hay un mecanismo de conmutación automática documentado.
Para aplicar CAP y PACELC haría falta conocer la replicación, el comportamiento
ante particiones y cómo se resuelve el compromiso entre latencia y consistencia
en operación normal. El informe permite identificar la dependencia de la base
de datos, pero deja esos puntos abiertos \cite[figura~4]{clab2026}.

\paragraph{Fallo observado durante las pruebas.}
En una prueba con Postman, un bloqueo por evento institucional dejó sin aviso
a los docentes con reservas afectadas. El equipo corrigió una condición en
\texttt{BloqueoService} durante el mismo sprint. Para verificar el flujo
completo, comprobaría que se registre el bloqueo y se genere el aviso \cite[sección~6.5.2]{clab2026}.

\paragraph{Propuesta y demostración.}
Para la resolución de fallas, propongo guardar el evento pendiente en la misma
transacción que actualiza el reporte. Así, si Mosquitto está caído, un proceso
puede reintentar la publicación y el aviso se puede demorar sin deshacer la
resolución. En MQTTX publico \texttt{ReporteFallaResuelto} en
\texttt{clab/fallas/resuelta} y muestro el mismo mensaje en Notificaciones y
Auditoría. La demo comprueba la entrega a dos suscriptores conectados con QoS~1.
El broker distribuye por tópico. Con QoS~1, la confirmación es por tramo y
pueden llegar duplicados. Cada consumidor podría reconocerlos mediante
\texttt{eventoId} \cite{mqtt2019}.

% si el archivo principal está en la raíz. Las citas usan referencias.bib.

\subsection{Carlos Manotoa · MQTT aplicado a CLAB}

\paragraph{Arquitectura y comunicación.}
CLAB gestiona los laboratorios informáticos de la UTEQ. Cuando el
administrador marca una falla como resuelta, Angular envía una petición REST a
Spring Boot y espera su respuesta. El backend actualiza el reporte en
PostgreSQL, guarda el seguimiento y avisa al docente que reportó la falla.
Para el administrador, todo empieza con un botón. La aplicación se encarga del
acceso a la base de datos \cite[sección~3.2 y caso CU-04]{clab2026}.

\paragraph{Despliegue y CAP/PACELC.}
El diagrama muestra un servidor PostgreSQL. Si Spring Boot pierde la conexión
con ese servidor, no hay un mecanismo de conmutación automática documentado.
Para aplicar CAP y PACELC haría falta conocer la replicación, el comportamiento
ante particiones y cómo se resuelve el compromiso entre latencia y consistencia
en operación normal. El informe permite identificar la dependencia de la base
de datos, pero deja esos puntos abiertos \cite[figura~4]{clab2026}.

\paragraph{Fallo observado durante las pruebas.}
En una prueba con Postman, un bloqueo por evento institucional dejó sin aviso
a los docentes con reservas afectadas. El equipo corrigió una condición en
\texttt{BloqueoService} durante el mismo sprint. Para verificar el flujo
completo, comprobaría que se registre el bloqueo y se genere el aviso \cite[sección~6.5.2]{clab2026}.

\paragraph{Propuesta y demostración.}
Para la resolución de fallas, propongo guardar el evento pendiente en la misma
transacción que actualiza el reporte. Así, si Mosquitto está caído, un proceso
puede reintentar la publicación y el aviso se puede demorar sin deshacer la
resolución. En MQTTX publico \texttt{ReporteFallaResuelto} en
\texttt{clab/fallas/resuelta} y muestro el mismo mensaje en Notificaciones y
Auditoría. La demo comprueba la entrega a dos suscriptores conectados con QoS~1.
El broker distribuye por tópico. Con QoS~1, la confirmación es por tramo y
pueden llegar duplicados. Cada consumidor podría reconocerlos mediante
\texttt{eventoId} \cite{mqtt2019}.

% si el archivo principal está en la raíz. Las citas usan referencias.bib.

\subsection{Carlos Manotoa · MQTT aplicado a CLAB}

\paragraph{Arquitectura y comunicación.}
CLAB gestiona los laboratorios informáticos de la UTEQ. Cuando el
administrador marca una falla como resuelta, Angular envía una petición REST a
Spring Boot y espera su respuesta. El backend actualiza el reporte en
PostgreSQL, guarda el seguimiento y avisa al docente que reportó la falla.
Para el administrador, todo empieza con un botón. La aplicación se encarga del
acceso a la base de datos \cite[sección~3.2 y caso CU-04]{clab2026}.

\paragraph{Despliegue y CAP/PACELC.}
El diagrama muestra un servidor PostgreSQL. Si Spring Boot pierde la conexión
con ese servidor, no hay un mecanismo de conmutación automática documentado.
Para aplicar CAP y PACELC haría falta conocer la replicación, el comportamiento
ante particiones y cómo se resuelve el compromiso entre latencia y consistencia
en operación normal. El informe permite identificar la dependencia de la base
de datos, pero deja esos puntos abiertos \cite[figura~4]{clab2026}.

\paragraph{Fallo observado durante las pruebas.}
En una prueba con Postman, un bloqueo por evento institucional dejó sin aviso
a los docentes con reservas afectadas. El equipo corrigió una condición en
\texttt{BloqueoService} durante el mismo sprint. Para verificar el flujo
completo, comprobaría que se registre el bloqueo y se genere el aviso \cite[sección~6.5.2]{clab2026}.

\paragraph{Propuesta y demostración.}
Para la resolución de fallas, propongo guardar el evento pendiente en la misma
transacción que actualiza el reporte. Así, si Mosquitto está caído, un proceso
puede reintentar la publicación y el aviso se puede demorar sin deshacer la
resolución. En MQTTX publico \texttt{ReporteFallaResuelto} en
\texttt{clab/fallas/resuelta} y muestro el mismo mensaje en Notificaciones y
Auditoría. La demo comprueba la entrega a dos suscriptores conectados con QoS~1.
El broker distribuye por tópico. Con QoS~1, la confirmación es por tramo y
pueden llegar duplicados. Cada consumidor podría reconocerlos mediante
\texttt{eventoId} \cite{mqtt2019}.
